\documentclass[10pt, aspectratio=169]{beamer}
% Preámbulo modular para presentaciones Beamer (Modelación Estadística)
% Diseño y paleta institucional del Tecnológico de Monterrey

\documentclass[aspectratio=169,xcolor={dvipsnames,table}]{beamer}

\usepackage[utf8]{inputenc}
\usepackage[spanish,mexico]{babel}
\usepackage{amsmath,amssymb,amsfonts,mathtools}
\usepackage{graphicx}
\usepackage{listings}
\usepackage{booktabs}
\usepackage{tikz}
\usetikzlibrary{matrix,arrows,decorations.pathmorphing,shapes.geometric,positioning}

% Paleta Institucional Tecnológico de Monterrey
\definecolor{TecAzul}{HTML}{0039A6}       % Azul TEC: color primario institucional
\definecolor{TecAzulOscuro}{HTML}{1A2E51} % Azul marino oscuro
\definecolor{TecRojo}{HTML}{EC2661}       % Rojo de acento (paleta miTec)
\definecolor{TecGris}{HTML}{646464}       % Gris neutro para comentarios y subtítulos
\definecolor{TecFondoBloque}{HTML}{F4F6F9}% Fondo suave para bloques

% Configuración de Tema Beamer
\usetheme{Boadilla}
\usecolortheme[named=TecAzulOscuro]{structure}
\setbeamercolor{alerted text}{fg=TecRojo}
\setbeamercolor{palette primary}{bg=TecAzulOscuro,fg=white}
\setbeamercolor{palette secondary}{bg=TecAzul,fg=white}
\setbeamercolor{palette tertiary}{bg=TecAzulOscuro!80!black,fg=white}
\setbeamercolor{title}{fg=TecAzulOscuro,bg=TecFondoBloque}
\setbeamercolor{frametitle}{fg=TecAzulOscuro,bg=TecFondoBloque}

% Bloques personalizados elegantes
\setbeamercolor{block title}{bg=TecAzulOscuro,fg=white}
\setbeamercolor{block body}{bg=TecFondoBloque,fg=black}
\setbeamercolor{block title alerted}{bg=TecRojo,fg=white}
\setbeamercolor{block body alerted}{bg=TecRojo!8,fg=black}
\setbeamercolor{block title example}{bg=TecAzul,fg=white}
\setbeamercolor{block body example}{bg=TecAzul!8,fg=black}

% Redondear bordes y añadir sombra sutil
\setbeamertemplate{blocks}[rounded][shadow=true]
\setbeamertemplate{navigation symbols}{}

% Pie de página limpio con número de diapositiva
\setbeamertemplate{footline}{
  \leavevmode%
  \hbox{%
  \begin{beamercolorbox}[wd=.7\paperwidth,ht=2.25ex,dp=1ex,left]{author in head/foot}%
    \hspace*{2ex}\insertshortauthor~~(\insertshortinstitute) \hfill \insertshorttitle
  \end{beamercolorbox}%
  \begin{beamercolorbox}[wd=.3\paperwidth,ht=2.25ex,dp=1ex,right]{date in head/foot}%
    \insertshortdate{}\hspace*{2em}
    \insertframenumber{} / \inserttotalframenumber\hspace*{2ex}
  \end{beamercolorbox}}%
  \vskip0pt%
}

% Configuración de Listings de Python para visualización pedagógica de código
\lstset{
    language=Python,
    basicstyle=\ttfamily\footnotesize,
    keywordstyle=\color{TecAzulOscuro}\bfseries,
    stringstyle=\color{TecRojo},
    commentstyle=\color{TecGris}\itshape,
    showstringspaces=false,
    breaklines=true,
    frame=lines,
    rulecolor=\color{TecAzulOscuro!30},
    backgroundcolor=\color{TecFondoBloque},
    aboveskip=0.5em,
    belowskip=0.5em,
    literate={á}{{\'a}}1 {é}{{\'e}}1 {í}{{\'i}}1 {ó}{{\'o}}1 {ú}{{\'u}}1
             {Á}{{\'A}}1 {É}{{\'E}}1 {Í}{{\'I}}1 {Ó}{{\'O}}1 {Ú}{{\'U}}1
             {ñ}{{\~n}}1 {Ñ}{{\~N}}1 {¿}{{?`}}1 {¡}{{!`}}1
}

% Metadatos institucionales por defecto
\title[Modelación Estadística]{Modelación Estadística para Ciencia de Datos e Ingeniería}
\author[J. Castillo Colmenares]{Juliho David Castillo Colmenares}
\institute[Tec de Monterrey]{Tecnológico de Monterrey \\ Escuela de Ingeniería y Ciencias}
\date{\today}

% Comandos y macros estadísticas compartidas para las presentaciones Beamer (Metropolis)

% Conjuntos numéricos básicos
\newcommand{\R}{\mathbb{R}}
\newcommand{\N}{\mathbb{N}}
\newcommand{\Q}{\mathbb{Q}}
\newcommand{\Z}{\mathbb{Z}}
\newcommand{\set}[1]{\left\{ #1 \right\}}
\newcommand{\abs}[1]{\left|#1\right|}
\newcommand{\norm}[1]{\left\| #1 \right\|}

% Comandos estadísticos y probabilísticos del curso
\newcommand{\Var}{\operatorname{Var}}
\newcommand{\cov}{\operatorname{Cov}}
\newcommand{\corr}{\rho}
\newcommand{\comb}[2]{\begin{pmatrix} #1 \\ #2 \end{pmatrix}}
\newcommand{\s}{\sigma}
\newcommand{\particion}{\mathfrak{P}}
\newcommand{\card}[1]{\#\left( #1 \right)}
\newcommand{\seq}[3]{\set{#1}_{#2}^{#3}}
\newcommand{\E}{\mathbb{E}}
\newcommand{\Prob}{\mathbb{P}}

% Entornos pedagógicos y cajas en español académico natural
\newenvironment{discusion}{%
  \begin{alertblock}{Pregunta de Discusión}%
}{%
  \end{alertblock}%
}

\newenvironment{reflexion}{%
  \begin{alertblock}{Reflexión para el Aula}%
}{%
  \end{alertblock}%
}

\newenvironment{paradoja}{%
  \begin{alertblock}{Paradoja de Exploración}%
}{%
  \end{alertblock}%
}

\newenvironment{motivacion}{%
  \begin{exampleblock}{Motivación: Ciencia de Datos en la Práctica}%
}{%
  \end{exampleblock}%
}

\newenvironment{retoclase}{%
  \begin{block}{Reto Rápido en Equipo}%
}{%
  \end{block}%
}


\title[Distribución Normal]{Sección 04.05: Distribución Normal}
\subtitle{Campana de Gauss, Estandarización al Puntaje $Z$ y Teorema del Límite Central}
\author[J. Castillo Colmenares]{Juliho Castillo Colmenares}
\institute{Tecnológico de Monterrey}
\date{\vspace{-1.2cm}}

\begin{document}

% Slide 1: Portada
\begin{frame}[plain]
  \vspace{-0.8cm}
  \titlepage
\end{frame}

% Slide 2: Hoja de Ruta
\begin{frame}{Hoja de Ruta --- Capítulo 04: Variables Aleatorias Continuas}
  \footnotesize
  ¿Dónde nos encontramos en el desarrollo modular del capítulo? \pause
  \begin{itemize}\itemsep=0.08em
    \item \textbf{04.01} Función de Densidad (PDF) y Soporte Continuo \pause
    \item \textbf{04.02} Función de Distribución Acumulada (CDF) Continua y Cuantiles \pause
    \item \textbf{04.03} Esperanza, Varianza y LOTUS Continuo \pause
    \item \textbf{04.04} Distribución Uniforme Continua $U(a, b)$ \pause
    \item \textbf{04.05} Distribución Normal / Gaussiana y Puntaje $Z$ \emph{(Hoy)} \pause
    \item \textbf{04.06} Distribución Exponencial y Procesos sin Memoria \pause
    \item \textbf{04.07} Distribuciones Gamma, Beta y Weibull
  \end{itemize}
  \vspace{-0.05cm}
  \begin{block}{Objetivo de la Sesión}
    Introducir la distribución Normal como la más importante de la estadística continua, dominar la estandarización al puntaje $Z$ y la regla 68-95-99.7, aplicar la MGF para demostrar propiedades de suma, y establecer el Teorema del Límite Central como justificación teórica de la universalidad de la Normal.
  \end{block}
\end{frame}

% Slide 3: Motivación
\begin{frame}{La Campana de Gauss: Universalidad de la Normal}
  \footnotesize
  \begin{motivacion}[¿Por qué la distribución Normal?]
    La distribución Normal $N(\mu, \sigma^{2})$ es sin duda la más importante de toda la estadística. Su forma de campana modela errores de medición, ruido en señales, alturas de personas, calificaciones de exámenes, y muchos otros fenómenos naturales. Por el Teorema del Límite Central, la suma (o promedio) de variables i.i.d. converge a la Normal independientemente de la distribución original.
  \end{motivacion}

  \vspace{0.15cm}
  \pause
  \begin{itemize}\itemsep=0.1cm
    \item \textbf{Forma de campana:} Simétrica alrededor de $\mu$ con dispersión controlada por $\sigma$. \pause
    \item \textbf{Estandarización:} Cualquier Normal se reduce a $N(0, 1)$ vía $Z = (X - \mu)/\sigma$. \pause
    \item \textbf{Universalidad:} Aparece en el TLC como límite de sumas; base de la inferencia estadística.
  \end{itemize}
\end{frame}

% Slide 4: Definición Formal
\begin{frame}{Definición Formal de la Distribución Normal}
  \footnotesize
  Una variable aleatoria continua $X$ tiene distribución \emph{Normal} con parámetros $\mu \in \mathbb{R}$ y $\sigma > 0$ si su PDF es: \pause

  \vspace{0.1cm}
  \begin{block}{Función de Densidad de Probabilidad (Forma de Campana)}
    $$f(x) = \dfrac{1}{\sigma\sqrt{2\pi}} \exp\left(-\dfrac{(x - \mu)^{2}}{2\sigma^{2}}\right), \quad x \in \mathbb{R}$$
    Por normalización: $\int_{-\infty}^{\infty} f(x)\,dx = 1$ (integral gaussiana). En este caso escribimos $X \sim N(\mu, \sigma^{2})$.
  \end{block}

  \vspace{0.1cm}
  \pause
  \begin{alertblock}{Características Fundamentales}
    \begin{itemize}\itemsep=0.05cm
      \item \textbf{Simetría:} $f(\mu - x) = f(\mu + x)$; la campana es perfectamente simétrica alrededor de $\mu$.
      \item \textbf{Moda, mediana, media:} Todas coinciden en $\mu$ (distribución unimodal).
      \item \textbf{Soporte:} Todo $\mathbb{R}$ (no acotada, dos colas infinitas).
    \end{itemize}
  \end{alertblock}
\end{frame}

% Slide 5: Estandarización
\begin{frame}{Estandarización: $Z = (X - \mu)/\sigma \sim N(0, 1)$}
  \footnotesize
  La \emph{estandarización} transforma cualquier Normal en la Normal estándar: \pause

  \vspace{0.1cm}
  \begin{block}{Puntaje Z (Z-Score)}
    $$Z = \dfrac{X - \mu}{\sigma} \sim N(0, 1)$$
    La inversa es $X = \mu + \sigma Z$. La estandarización preserva las probabilidades: $P(a \le X \le b) = P\left(\frac{a-\mu}{\sigma} \le Z \le \frac{b-\mu}{\sigma}\right)$.
  \end{block}

  \vspace{0.1cm}
  \pause
  \begin{alertblock}{CDF Normal Estándar y Simetría}
    \begin{itemize}\itemsep=0.05cm
      \item $\Phi(z) = P(Z \le z) = \int_{-\infty}^{z} \dfrac{1}{\sqrt{2\pi}} e^{-u^{2}/2}\,du$ (CDF estándar).
      \item $\Phi(-z) = 1 - \Phi(z)$ (simetría de la campana).
      \item $\Phi(0) = 0.5$, $\Phi(1) \approx 0.8413$, $\Phi(2) \approx 0.9772$, $\Phi(3) \approx 0.9987$.
    \end{itemize}
  \end{alertblock}
\end{frame}

% Slide 6: Regla 68-95-99.7
\begin{frame}{Regla Empírica 68-95-99.7}
  \footnotesize
  La \emph{regla empírica} de la Normal proporciona una guía rápida para interpretar probabilidades: \pause

  \vspace{0.1cm}
  \begin{block}{Probabilidades dentro de $k$ Desviaciones Estándar}
    \begin{itemize}\itemsep=0.06cm
      \item $P(\mu - \sigma \le X \le \mu + \sigma) = 2\Phi(1) - 1 \approx 0.6827$ (68\%).
      \item $P(\mu - 2\sigma \le X \le \mu + 2\sigma) = 2\Phi(2) - 1 \approx 0.9545$ (95\%).
      \item $P(\mu - 3\sigma \le X \le \mu + 3\sigma) = 2\Phi(3) - 1 \approx 0.9973$ (99.7\%).
    \end{itemize}
  \end{block}

  \vspace{0.1cm}
  \pause
  \begin{alertblock}{Aplicación: Pruebas de Hipótesis}
    \begin{itemize}\itemsep=0.05cm
      \item Bajo $H_{0}$, un estadístico $Z \sim N(0, 1)$ cae fuera de $\pm 1.96$ solo el 5\% de las veces.
      \item Si $|Z_{\text{obs}}| > 1.96$, rechazamos $H_{0}$ a nivel $\alpha = 0.05$ (prueba bilateral).
      \item Para pruebas unilaterales: $z_{0.025} = 1.96$ y $z_{0.05} = 1.645$ son los valores críticos estándar.
    \end{itemize}
  \end{alertblock}
\end{frame}

% Slide 7: MGF y Suma de Normales
\begin{frame}{MGF y Suma de Variables Normales Independientes}
  \footnotesize
  La MGF de la Normal tiene forma cerrada simple, y la suma de normales independientes es Normal: \pause

  \vspace{0.1cm}
  \begin{block}{MGF Normal}
    $$M_{X}(t) = \E[e^{tX}] = \exp\left(\mu t + \dfrac{\sigma^{2} t^{2}}{2}\right)$$
    Derivando: $M'(0) = \mu$ y $M''(0) = \mu^{2} + \sigma^{2}$, confirmando $\E[X] = \mu$ y $\Var(X) = \sigma^{2}$.
  \end{block}

  \vspace{0.1cm}
  \pause
  \begin{alertblock}{Estabilidad bajo Suma}
    \begin{itemize}\itemsep=0.05cm
      \item Si $X \sim N(\mu_{1}, \sigma_{1}^{2})$ y $Y \sim N(\mu_{2}, \sigma_{2}^{2})$ independientes, entonces $X + Y \sim N(\mu_{1} + \mu_{2}, \sigma_{1}^{2} + \sigma_{2}^{2})$.
      \item La familia Normal es \emph{cerrada} bajo suma (e independientes): la suma de normales es Normal.
      \item Para $X_{1}, \dots, X_{n}$ i.i.d. $N(\mu, \sigma^{2})$: $\bar{X} \sim N(\mu, \sigma^{2}/n)$.
    \end{itemize}
  \end{alertblock}
\end{frame}

% Slide 8: Teorema del Límite Central
\begin{frame}{Teorema del Límite Central: Universalidad de la Normal}
  \footnotesize
  El \emph{Teorema del Límite Central} (TLC) establece que la suma de variables i.i.d. converge a la Normal: \pause

  \vspace{0.1cm}
  \begin{block}{Enunciado del TLC}
    Si $X_{1}, \dots, X_{n}$ son i.i.d. con media $\mu$ y varianza $\sigma^{2} < \infty$, entonces:
    $$Z_{n} = \dfrac{\bar{X} - \mu}{\sigma/\sqrt{n}} \xrightarrow{d} N(0, 1) \quad \text{cuando } n \to \infty$$
    La distribución límite es Normal \emph{universal}, independientemente de la distribución original de $X_{i}$.
  \end{block}

  \vspace{0.1cm}
  \pause
  \begin{alertblock}{Implicaciones Prácticas}
    \begin{itemize}\itemsep=0.05cm
      \item \textbf{Aproximación de la Binomial:} Para $n$ grande, $\text{Bin}(n, p) \approx N(np, np(1-p))$.
      \item \textbf{Convergencia rápida:} Para distribuciones simétricas, $n = 20$ es suficiente; para asimétricas, $n \ge 50$.
      \item \textbf{Base de pruebas estadísticas:} $Z = (\bar{X} - \mu)/(\sigma/\sqrt{n}) \sim N(0, 1)$ bajo $H_{0}$ para $n$ grande.
    \end{itemize}
  \end{alertblock}
\end{frame}

% Slide 9: Aplicaciones
\begin{frame}{Aplicaciones en Inferencia, Ciencias y Ciencia de Datos}
  \footnotesize
  La distribución Normal es la más usada en análisis cuantitativo: \pause

  \vspace{0.15cm}
  \begin{itemize}\itemsep=0.12cm
    \item \textbf{Inferencia Estadística:} Pruebas $Z$, intervalos de confianza, regresión lineal, y análisis de varianza asumen Normalidad. \pause
    \item \textbf{Calidad y Six Sigma:} Cartas de control $\bar{x}$ y $R$ basadas en $\pm 3\sigma$ de la Normal. \pause
    \item \textbf{Finanzas:} Modelo de Black-Scholes para precios de opciones; movimientos Brownianos geométricos. \pause
    \item \textbf{Machine Learning:} Regresión lineal Gaussiana, modelos lineales generalizados, y aproximación de funciones de verosimilitud.
  \end{itemize}
\end{frame}

% Slide 12A-1: Lab Python (Bloque 1)
\begin{frame}[fragile]{Laboratorio en Python: PDF y Estandarización}
  \scriptsize
  Validación de $\int f = 1$, estandarización al puntaje $Z$, y verificación de la regla 68-95-99.7 (\texttt{04.05\_normal\_distribution.py}):
  \vspace{0.02cm}
  {\lstset{basicstyle=\fontsize{4.5pt}{5.5pt}\selectfont\ttfamily}
  \lstinputlisting[firstline=15, lastline=33]{../../code/04_variables_aleatorias_continuas/04.05_normal_distribution.py}}
\end{frame}

% Slide 12A-2: Lab Python (Bloque 2)
\begin{frame}[fragile]{Laboratorio en Python: Prueba $Z$ y Cuantiles}
  \scriptsize
  Cálculo de probabilidades de intervalo, cuantiles y prueba de hipótesis $Z$ de una muestra:
  \vspace{0.02cm}
  {\lstset{basicstyle=\fontsize{4pt}{5pt}\selectfont\ttfamily}
  \lstinputlisting[firstline=40, lastline=63]{../../code/04_variables_aleatorias_continuas/04.05_normal_distribution.py}}
\end{frame}

% Slide 12B: Lab Python (Bloque 3)
\begin{frame}[fragile]{Laboratorio en Python: TLC y Aproximación Binomial-Normal}
  \scriptsize
  Verificación Monte Carlo del TLC y aplicación de la aproximación Binomial-Normal con corrección de Yates:
  \vspace{0.0cm}
  {\lstset{basicstyle=\fontsize{4pt}{5pt}\selectfont\ttfamily}
  \lstinputlisting[firstline=73, lastline=100]{../../code/04_variables_aleatorias_continuas/04.05_normal_distribution.py}}
\end{frame}

% Slide 12C: Salida Terminal
\begin{frame}[fragile]{Laboratorio en Python: Salida en Terminal}
  \scriptsize
  Salida estándar generada al ejecutar el script del laboratorio en Python:
  \vspace{0.02cm}
  \begin{block}{Salida Estándar en Consola}
    \fontsize{4pt}{5pt}\selectfont
    \begin{verbatim}
=== Block 1: Normal PDF & Standardization ===
Normal(100,15) PDF integral: 1.00000000
Z-score: F(x) actual = Phi(Z) (diff = 0)
Empirical rule: 68.27% | 95.45% | 99.73%

=== Block 2: Z-Score & Interval Probabilities ===
Heights N(170,100): P(X>185)=0.0668, P(160<=X<=180)=0.6827
Quantiles: q_0.025=-1.96, q_0.05=-1.64, q_0.5=0, q_0.95=1.64, q_0.975=1.96
Z-test: Z=1.5, p-value=0.1336, fail to reject H0

=== Block 3: CLT & Normal Approximation ===
Sum of 30 U(0,1) ~ N(15.0, 2.5): KS D=0.0014, p=0.71
Binomial(400, 0.25) approx:
  Exact: P(X>=120)=0.0133
  Normal w/ Yates: 0.0122 (error 0.0011)
  Required n for SE<0.01: 1875
    \end{verbatim}
  \end{block}
\end{frame}

% Slide 17: Cierre
\begin{frame}{Cierre de Sección y Síntesis sobre la Distribución Normal}
  \begin{reflexion}[La Campana de Gauss como Distribución Universal]
    La distribución Normal $N(\mu, \sigma^{2})$ es sin duda la más importante de toda la estadística. Su forma de campana simétrica, parametrizada por solo dos momentos, captura la esencia de muchos fenómenos naturales. El Teorema del Límite Central establece su universalidad: la suma de variables i.i.d. converge a la Normal independientemente de la distribución original.
  \end{reflexion}

  \vspace{0.2cm}
  \pause
  \begin{block}{Lecciones Clave de la Distribución Normal}
    \begin{itemize}\itemsep=0.08cm
      \item \textbf{PDF:} $f(x) = \frac{1}{\sigma\sqrt{2\pi}} e^{-(x-\mu)^{2}/(2\sigma^{2})}$ modela errores y fluctuaciones.
      \item \textbf{Estandarización:} $Z = (X-\mu)/\sigma$ reduce a $N(0,1)$; las probabilidades se leen de tablas universales.
      \item \textbf{TLC:} $Z_{n} = (\bar{X}-\mu)/(\sigma/\sqrt{n}) \xrightarrow{d} N(0,1)$ explica la omnipresencia de la Normal.
    \end{itemize}
  \end{block}
\end{frame}

% Slide 18: Perspectiva Modular
\begin{frame}{Perspectiva Modular --- El Próximo Reto}
  \scriptsize
  \begin{retoclase}[Para la próxima sesión: Distribuciones Gamma, Beta y Weibull]
    La Normal modela procesos con muchos sumandos pequeños. Ahora estudiaremos distribuciones que modelan \emph{tiempos de espera} con formas flexibles: la \emph{Gamma} generaliza la Exponencial y la Erlang con parámetro de forma $\alpha$ no necesariamente entero; la \emph{Beta} es la distribución canónica para variables en $[0, 1]$ (proporciones, probabilidades); y la \emph{Weibull} es fundamental en análisis de confiabilidad con tasa de fallo variable en el tiempo.
  \end{retoclase}

  \vspace{0.08cm}
  \pause
  \begin{alertblock}{Preparación Recomendada}
    \begin{itemize}\itemsep=0.06cm
      \item Repasa la función Gamma $\Gamma(\alpha) = \int_{0}^{\infty} t^{\alpha-1} e^{-t}\,dt$ y sus propiedades.
      \item Reflexiona sobre cómo la flexibilidad de forma de la Gamma permite modelar distribuciones con diferentes grados de asimetría.
    \end{itemize}
  \end{alertblock}
\end{frame}

\end{document}
